\chapter{Protokollok ismertetése}

\section{Kyber protocoll}

A CRYSTALS KYBER protokoll a posztkvantum-kriptográfia egyik jelentős megoldása, amely a NIST
által írt elemzések alapján biztonságosak. Ez a protokoll a rács alapú titkosításra épül, amit LWE
alapokon valósít meg, és számos biztonsági tulajdonsággal rendelkezik, amelyek lehetővé teszik a
megbízható és titkos kommunikációt.

A KYBER protokoll alapvetően három fő részből áll: a kulcskészítés és kulcscsere folyamatából, az
üzenet titkosításából és az üzenet visszafejtéséből.

\subsection{Paraméterek és kulcskészítés}

A kulcskészítés folyamatának első lépése a paraméterek inicializálása. Ez magában foglalja a rács
paramétereinek és a biztonsági szintnek a meghatározását. A KYBER protokoll különböző biztonsági
szinteket kínál, amelyek lehetővé teszik a felhasználók számára a megfelelő szintű védelem kiválasztását.

A kulcskészítés során a Alice (a fogadó fél) elkészíti a privát és publikus kulcsokat. A privát kulcs lesz az
s vektor, a publikus pedig az A mátrix és a t vektor. Ez így az LWE séma kulcsgenerálási folyamata

\subsection{Üzenet titkosítása}
A Bob (a küldő fél) az Alice publikus kulcsát használja az üzenet titkosításához. A Bob egy véletlenszerű vektort generál, majd a publikus kulcsot használva létrehozza a titkosított üzenetet. Ez a folyamat magában foglalja a rács alapú műveleteket és a hibák hozzáadását, hogy biztosítsa a titkosítás biztonságát.

Az üzenet titkosításakor Bob (a küldő fél) egy saját titkos kulcs segítségével és újabb hibavektorokkal
titkosítja üzenetét, amely az u és v vektorok lesznek.

\subsection{Titkosítás visszafejtése}

Visszafejtéskor Alice a titkos kulcsát felhasználva megkapja az üzenetet és egy hibavektort, amit a bázis
ismeretében a megfelelő értékre le tudja kerekíteni.

Az algoritmus hatékonysága és teljesítménye azt jelenti, hogy a KYBER protokoll képes-e a valós
környezetben megfelelő sebességgel és erőforrásokkal működni. Ez kritikus fontosságú a gyakorlati
alkalmazások szempontjából, mivel az információvédelem mellett a protokoll gyors és hatékony
működését is biztosítania kell.


\section{Kriptográfiai protokollok}
Először a Kyber protokollt mutatjuk be, amely egy posztkvantum kriptográfiai megoldás, és a dolgozatban az IoT rendszerek biztonságának megerősítésére alkalmazzuk. Részletesen ismertetjük a Kyber működését, beleértve a kulcscserét és az üzenet titkosítását, valamint a protokoll biztonsági tulajdonságait. Ezután bemutatjuk a BIKE protokollt, amely szintén egy posztkvantum kriptográfiai megoldás, és a dolgozatban a Kyber protokoll alternatívájaként szerepel. Ismertetjük a BIKE működését, beleértve a kulcscserét és az üzenet titkosítását, valamint a protokoll biztonsági tulajdonságait. Végül összehasonlítjuk a Kyber és a BIKE protokollokat, kiemelve előnyeiket és hátrányaikat az IoT rendszerek biztonságának megerősítése szempontjából.

\section{CRYSTALS-Kyber protokoll}

\section{BIKE protokoll}